\section{Experiments}
\label{sec:experiments}

We evaluate the proposed framework on two complementary problems:
(i) a controlled 2D toy task that allows us to verify whether the
recovered routing matches the ground-truth cluster structure, and
(ii) a heterogeneous CIFAR-100 / SVHN image mixture, where the data
source defines an oracle partition and we compare against several MoE
and NAS baselines.

\subsection{Toy Experiment}
\label{subsec:toy}

\paragraph{Setup.}
We sample $2000$ points in $\mathbb{R}^2$ from a four-class problem
consisting of one ``linear'' cloud (two classes separated by a
diagonal) and one ``ring'' cloud (two classes corresponding to an
inner and an outer ring). The data is split $80/20$ into
training/validation and clustered into $M = 20$ KMeans clusters,
yielding the partition described in our data preparation pipeline.
We use $K = 2$ experts so that the ideal solution is unambiguous: one
expert should specialise to the linear cloud and the other to the
ring cloud. The search space is the cell-based DARTS-style space of
\cite{liu2018darts}, restricted to small operations suitable for 2D
inputs.

\paragraph{Results.}
With the parameters reported in our experimental log, SAGEM
(Algorithm~\ref{alg:sgem}) recovers the \emph{ideal} cluster split
$(10, 10)$ between the two experts in $20$ EM iterations, matching
the hand-defined oracle assignment that splits clusters according to
the centroid coordinate $x_0 = -2$ (Figure~\ref{fig:toy}). This sanity check confirms that
the surrogate-guided objective in~\eqref{eq:moe_objective} is well
posed and that the proposed optimiser solves it when the surrogate
is sufficiently accurate, in agreement with the convergence
guarantee of Theorem~\ref{thm:sgem}.

\begin{figure}[h]
    \centering
    \begin{subfigure}[b]{0.32\textwidth}
        \centering
        \includegraphics[width=\textwidth]{toy_dataset}
        \caption{}
        \label{fig:toy_dataset}
    \end{subfigure}
    \hfill
    \begin{subfigure}[b]{0.32\textwidth}
        \centering
        \includegraphics[width=\textwidth]{toy_clusters}
        \caption{}
        \label{fig:toy_clusters}
    \end{subfigure}
    \hfill
    \begin{subfigure}[b]{0.32\textwidth}
        \centering
        \includegraphics[width=\textwidth]{toy_sgem_split}
        \caption{}
        \label{fig:toy_sgem}
    \end{subfigure}
    \caption{Toy experiment. (a) The four-class dataset: a linear cloud
    (classes 0--1) and a ring cloud (classes 2--3). (b) The $M = 20$
    KMeans clusters (crosses mark the centroids). (c) The expert
    assignment recovered by SAGEM with $K = 2$: it splits the clusters
    exactly along the linear/ring boundary, matching the oracle
    partition up to a permutation of expert labels.}
    \label{fig:toy}
\end{figure}

\subsection{CIFAR-100 / SVHN mixture}
\label{subsec:cifar}

\paragraph{Setup.}
To test whether the cluster-aware objective recovers a meaningful
data partition on realistic data, we build a heterogeneous benchmark
by mixing CIFAR-100 and SVHN in equal proportions ($42\,000$ images
total). SVHN labels are shifted by $100$, so the combined problem has
$110$ classes. The two sources form a maximally separated domain
structure, which gives a well-defined oracle: with $K = 2$ experts the
ideal routing assigns each cluster to the expert specialised on its
source. We partition the data into $M = 30$ \emph{semantic} clusters
obtained from a pretrained feature extractor, using a three-way
train/validation/test split ($0.7/0.15/0.15$). KMeans is fit on the
training embeddings only; validation and test points are assigned to
the nearest centroid. The resulting oracle partition is
$19$ CIFAR-majority and $11$ SVHN-majority clusters, with a mean
source purity of $0.978$.

All MoE configurations use $K = 2$ experts with
\texttt{init\_channels = 16}. The base network is a small DARTS-style
architecture \texttt{CIFAR100Net} (stem, two normal cells, one fixed
reduction cell, global average pooling, fully connected head). The MoE
wrapper instantiates $K$ such networks combined through a fixed hard
\emph{cluster gate}: each cluster is routed to a single expert, so the
gate directly reflects the discovered data partition. Architecture
search uses the train/validation split; the final model is retrained
for $100$ epochs on $\text{train} \cup \text{validation}$ and evaluated
\emph{once} on the held-out test set. We report test accuracy averaged
over five seeds. Training is performed inside a Docker container with
PyTorch~2.5.1 on a single RTX~3080~Ti GPU.

\paragraph{Baselines.}
We compare against the following methods, all evaluated with the same
cluster gate:
\begin{itemize}
    \item \textbf{Random-MoE}: $200$ random architecture pairs are
    searched at $30$ epochs each on a random cluster split, then the
    best is retrained.
    \item \textbf{DARTS$\times K$}: $K$ experts that share the same
    DARTS-discovered architecture on a random cluster split.
    \item \textbf{SAGEM}: our method with the Phase~C uniform-mixture
    coefficient $p_C = 0.5$ and load-balancing weight
    $\lambda_{\text{LB}} = 50$.
\end{itemize}

\paragraph{Main results.}
Table~\ref{tab:cifar_main} reports the test accuracy with cluster
gating. The differences between the methods are small and lie
within one standard deviation of each other: SAGEM
($0.5935 \pm 0.003$) is statistically indistinguishable from
DARTS$\times K$ ($0.5913 \pm 0.009$), and both improve only marginally
over the random-MoE baseline ($0.5551 \pm 0.031$, with a large
seed-to-seed variance). Hard cluster gating partitions the data so
that each expert sees fewer examples, which largely offsets the benefit
of specialisation, leaving the learned routing of SAGEM no better than
sharing a single architecture across experts on a random split.

\begin{table}[h]
\centering
\caption{Test accuracy on the CIFAR-100 / SVHN mixture with cluster
gating ($M=30$, $K=2$, semantic clustering, mean $\pm$ std over five
seeds).}
\label{tab:cifar_main}
\begin{tabular}{lc}
\toprule
Method & test acc.\ (cluster gate) \\
\midrule
Random-MoE & $0.5551 \pm 0.031$ \\
DARTS$\times K$ (same arch) & $0.5913 \pm 0.009$ \\
SAGEM, $\lambda_{\text{LB}}{=}50$, $p_C{=}0.5$ & $\mathbf{0.5935 \pm 0.003}$ \\
\bottomrule
\end{tabular}
\end{table}

\begin{figure}[p]
    \centering
    \includegraphics[height=0.92\textheight]{svhn_cluster_reps}
    \caption{Per-cluster representatives of the CIFAR-100 / SVHN
    mixture under the SAGEM routing ($K=2$, seed $322$; first $15$ of
    the $M = 30$ semantic clusters shown). Each row is one cluster: the
    leftmost column is the cluster medoid and the rest are the eight
    images nearest to the cluster mean. The medoid border (leftmost
    column) is coloured by the expert assigned by SAGEM (E0 red,
    E1 green).}
    \label{fig:svhn_reps}
\end{figure}

\paragraph{Does SAGEM recover the data source?}
Because the oracle partition by data source is sharp (mean purity
$0.978$), we can directly check whether the learned routing aligns
with it (Figure~\ref{fig:svhn_reps}). Across five seeds, the SAGEM
hard assignment agrees with the oracle source split on only
$\approx 60\%$ of clusters (best matching of the two label
permutations), barely above the $\approx 54\%$ baseline of a balanced
random split. The load-balancing penalty is the main culprit: the
oracle split is imbalanced ($19/11$) whereas the penalty pulls the
routing towards $15/15$, washing out the source signal. Hence the
hypothesis that the cluster-aware objective discovers the underlying
data domains is \emph{not} supported on this benchmark.

\paragraph{Discussion.}
With cluster gating, all MoE variants cluster tightly around
$0.59$ test accuracy, and the source-recovery analysis shows
that the discovered partition is essentially uncorrelated with the
true data structure. We attribute this to a self-reinforcing bias in
the surrogate dataset rather than to the EM optimiser itself: the
active-learning loop over-samples cluster masks favourable to a
dominant expert, so the routing is driven by the surrogate's bias and
is held balanced only by the load-balancing penalty. Removing this
bias --- for instance by symmetrising the data collection across
experts or enforcing permutation invariance of the surrogate --- is
the most promising direction for making the cluster-aware routing
informative. As Figure~\ref{fig:svhn_reps} illustrates at the level of
individual images, many SVHN-dominated clusters (digit crops) and
CIFAR-dominated clusters (natural images) are routed to different
experts, but the assignment is far from a clean source split,
consistent with the $\approx 60\%$ agreement reported above.
